\chapter{Bourses, courtiers, fournisseurs et régulateurs comme employeurs}\label{ch:in:exchanges-brokers-vendors-and-regulators-as-employers}

Les groupes de bourses cotés publient ce que les sociétés de trading privées ne publient pas~: leur revenu, leur résultat opérationnel et le nombre
de personnes qu'elles emploient, chaque année. En 2025, les quatre plus grands groupes américains employaient entre 1\,661 et 12\,844 personnes chacun et
dégageaient entre 0.24~million de dollars et 1.09~million de dollars de résultat opérationnel par employé. Le régulateur du marché publie
davantage encore~: la \omterm{def:in:exchanges-brokers-vendors-and-regulators-as-employers:band}{bande de rémunération} de chaque échelon de son personnel et le supplément propre à chaque ville, dans la limite d'un plafond de 292\,300~dollars en 2026.
Le reste de l'écosystème --- bourses, chambres de compensation, courtiers, fournisseurs de données et de technologie, régulateurs --- est
la partie de l'industrie la plus transparente pour y travailler, et celle dont les emplois sont les plus proches de l'ingénierie et du service public. Ce
chapitre le décrit en tant qu'employeur.

\section{Groupes de bourses et chambres de compensation}

\begin{definition}[Infrastructure de marché financier]\label{def:in:exchanges-brokers-vendors-and-regulators-as-employers:fmi}
Une \emph{infrastructure de marché financier}\index{infrastructure de marché financier} est un système multilatéral entre ses
participants pour la compensation, le règlement-livraison ou l'enregistrement de paiements, de titres, de dérivés ou d'autres transactions~: un
système de paiement d'importance systémique, un dépositaire central de titres, un système de règlement-livraison de titres, une
contrepartie centrale ou un référentiel de transactions. Les bourses et les plateformes de négociation sont les marchés que ces systèmes desservent, et elles appartiennent souvent
aux mêmes groupes.
\end{definition}

Un groupe de bourses vend des services de négociation, de compensation, de données, des licences d'indices et de connectivité (livre~1, chapitre~4~; livre~14). C'est une
entreprise technologique réglementée comme un service public~: ses ingénieurs construisent et exploitent des moteurs d'appariement (matching engine), des systèmes de données de marché et
des plateformes de compensation (livre~10, chapitre~27, en construit un simulateur)~; ses quants conçoivent des modèles de marge, des méthodologies d'indices
et de surveillance~; ses équipes produit conçoivent les contrats et les grilles tarifaires. Le travail se situe de l'autre côté des
problèmes des sociétés de trading~: la latence qu'elles se disputent est celle qu'un ingénieur d'une bourse mesure et rend équitable.

\begin{dated}{2025-12}{Quatre groupes de bourses dans leurs déclarations}\label{dat:in:exchanges-brokers-vendors-and-regulators-as-employers:exchanges}
\footnotesize
\begin{tabular}{@{}lrrrrr@{}}
\toprule
groupe (10-K 2025) & revenu, \$m & rés.\ op., \$m & effectifs & revenu/tête & rés.\ op./tête\\
\midrule
CME Group & 6\,520.6 & 4\,229.5 & environ 3\,875 & 1.68~M\$ & 1.09~M\$\\
Cboe Global Markets & 4\,714.2 & 1\,467.1 & 1\,661 & 2.84~M\$ & 0.88~M\$\\
Intercontinental Exchange & 12\,640 & 4\,929 & 12\,844 & 0.98~M\$ & 0.38~M\$\\
Nasdaq & 8\,262 & 2\,331 & 9\,525 & 0.87~M\$ & 0.24~M\$\\
\bottomrule
\end{tabular}

\smallskip\noindent Le revenu est le total publié par chaque groupe, qui pour trois d'entre eux inclut les charges liées aux
transactions (remises et frais refacturés aux participants)~; le résultat opérationnel est calculé après elles. Deux des groupes publient des segments hors
négociation~: l'un publie les bourses, les titres à revenu fixe et les services de données, ainsi que la technologie hypothécaire~; un autre
publie les plateformes d'accès au capital, de technologie financière et de services de marché.
\end{dated}

Les quatre groupes se distinguent davantage par leur résultat opérationnel par tête que par tout autre élément
(voir la \cref{fig:in:exchanges-brokers-vendors-and-regulators-as-employers:perhead})~: le groupe dont le revenu vient surtout de la négociation et de la compensation de dérivés gagne
1.09~million de dollars par employé, soit 4.5 fois
celui du groupe qui publie un segment de technologie financière.
Le revenu par tête est moins comparable, car les frais refacturés (pass-through fees) gonflent le revenu des groupes qui versent des remises. Pour un
employé, la différence tient au métier~: une bourse de dérivés emploie des personnes proches de la négociation et de la compensation, un
groupe diversifié en emploie bien davantage dans le logiciel, les données et les services.

\begin{omfigure}
\begin{tikzpicture}
\begin{axis}[width=10.5cm,height=5.2cm,xbar,area legend,bar width=7pt,xmin=0,xmax=3.2,ymin=-0.6,ymax=3.6,xlabel={\small millions de dollars par employé, 2025},
  ytick={0,1,2,3},yticklabels={{CME Group},{Cboe Global Markets},{Intercontinental Exchange},{Nasdaq}},
  y dir=reverse,tick label style={font=\footnotesize},grid=major,grid style={gray!20},table/col sep=comma,
  legend style={legend cell align=left,font=\scriptsize,at={(0.5,-0.3)},anchor=north,/tikz/every even column/.append style={column sep=8pt}},legend columns=2]
\addplot[fill=omDef!55,draw=omDef] table[y=k,x=opincome]{figdata/industry/09-exchanges-brokers-vendors-and-regulators-as-employers/perhead.csv};
\addplot[fill=gray!40,draw=gray] table[y=k,x=revenue]{figdata/industry/09-exchanges-brokers-vendors-and-regulators-as-employers/perhead.csv};
\legend{résultat opérationnel par tête,revenu par tête}
\end{axis}
\end{tikzpicture}

\omcaption{Résultat opérationnel et revenu par employé de quatre groupes de bourses américains, 2025, classés par résultat opérationnel par
tête. Le revenu inclut les charges liées aux transactions pour trois des groupes. Données~: formulaires 10-K pour 2025 et SEC XBRL
company facts, par l'intermédiaire de \texttt{in\_infra.exchanges}.}\label{fig:in:exchanges-brokers-vendors-and-regulators-as-employers:perhead}
\end{omfigure}

\section{Courtiers et fournisseurs de technologie de trading}

Les courtiers vendent l'accès et l'exécution~: les courtiers d'agence acheminent et travaillent les ordres (livre~10, chapitre~20), les prime brokers financent et
la conservation (livre~1, chapitre~6), les courtiers compensateurs se placent entre leurs clients et les chambres de compensation (livre~1, chapitre~30).
Les fournisseurs de technologie de trading vendent les systèmes que les sociétés ne construisent pas elles-mêmes~: systèmes de gestion des ordres et systèmes de gestion de l'exécution
(EMS), flux de données de marché et ticker plants, logiciels de risque et de conformité, connectivité (livre~14, chapitre~27, en dresse
la carte). Pour un quant ou un ingénieur, ces employeurs offrent les mêmes problèmes techniques qu'une société de trading --- latence,
exactitude, passage à l'échelle --- avec pour enjeu le P\&L des clients de la société plutôt que le sien, et une rémunération
qui suit l'économie d'une entreprise de logiciels~: salaires et actions plutôt qu'une part du profit de trading.

\section{Fournisseurs de données}

Les fournisseurs de données de marché collectent, normalisent, stockent et redistribuent des données~: flux des bourses, données de référence, cours des
instruments négociés de gré à gré, données alternatives (livre~7, chapitre~12). Leurs emplois quantitatifs portent sur la qualité des données
évaluation des prix (modèles qui valorisent les obligations et les dérivés n'ayant pas été négociés), construction d'indices et
analyses. Leurs emplois d'ingénierie sont ceux d'une plateforme de données à grande échelle (livre~15). Le travail est le plus proche de ce que fait une
équipe plateforme de recherche au sein d'une société de trading, mais pour de nombreux clients à la fois.

\section{Régulateurs et banques centrales}

Les régulateurs emploient des économistes, des analystes quantitatifs, des examinateurs, des juristes et des technologues. Leurs équipes de structure de marché et
de surveillance utilisent les méthodes de cette série sur l'ensemble du marché~: l'étude de la course à la latence du chapitre~2 relève de la
recherche des régulateurs, et la surveillance s'appuie sur les détecteurs du livre~9 (chapitre~29) sur le consolidated audit trail. Les banques centrales
emploient des quants dans les opérations de marché, la gestion des réserves et l'analyse de la stabilité financière. La rémunération est publique.

\begin{definition}[Bande de rémunération]\label{def:in:exchanges-brokers-vendors-and-regulators-as-employers:band}
Une \emph{bande de rémunération}\index{bande de rémunération} est la plage publiée du salaire de base pour un échelon de la grille salariale d'un employeur, allant
d'un minimum à un maximum~; la grille d'un employeur public publie aussi les suppléments (par exemple par ville) qui s'y appliquent
et tout plafond sur la rémunération totale.
\end{definition}

\begin{dated}{2026-04}{La grille salariale d'un régulateur}\label{dat:in:exchanges-brokers-vendors-and-regulators-as-employers:sec}
Le salaire de base 2026 de la SEC pour sa grille SK, en vigueur au 5 avril 2026, va de 24\,728 à 36\,213~dollars à l'échelon 1 jusqu'à
138\,383 à 239\,449~dollars à l'échelon 17 (superviseur). Échelons retenus~: SK-12 83\,718 à 141\,843~dollars~; SK-13 99\,557 à 168\,668~dollars~;
SK-14 114\,315 à 193\,670~dollars~; SK-16 130\,550 à 221\,711~dollars. Les taux de majoration locale s'appliquent au salaire de base selon le lieu d'affectation~: New York
37.95~\%, San Francisco 46.34~\%, Washington DC 33.94~\%, Chicago 30.86~\%. «~Le plafond salarial, majoration locale comprise, est
292\,300~dollars pour l'ensemble du personnel SK.~»
\end{dated}

\begin{method}[Une bande publique comme plage annuelle]\label{met:in:exchanges-brokers-vendors-and-regulators-as-employers:band}
\begin{enumerate}
\item Prendre le minimum et le maximum du salaire de base de l'échelon.
\item Multiplier les deux par un plus le taux de majoration locale du lieu d'affectation.
\item Plafonner chaque borne au plafond de la grille.
\item Comparer la plage obtenue à celle qu'une autre source donne pour la même profession (chapitre~14), en mesurant la
part de l'autre plage que couvre la bande.
\end{enumerate}
\end{method}

À New York, un poste SK-14 est rémunéré de 157\,698~\$ à 267\,168~\$~; un poste SK-16 va de 180\,094~\$ à 292\,300~\$, la borne haute
plafonnée (\cref{fig:in:exchanges-brokers-vendors-and-regulators-as-employers:bands}). Les bandes sont larges --- le haut de l'échelon
vaut environ 1.7 fois son bas --- car elles couvrent la progression d'une carrière au sein de l'échelon, et une nouvelle
recrue est en général placée près du bas. Le plafond joue sur les échelons supérieurs~: à partir du SK-16, la borne haute à New York est le plafond, et non
l'échelon.

\begin{omfigure}
\begin{tikzpicture}
\begin{axis}[width=10.5cm,height=5.6cm,xmin=8.4,xmax=17.6,ymin=0,ymax=340,xtick={9,10,11,12,13,14,15,16,17},
  xlabel={\small Échelon SK},ylabel={\small salaire annuel à New York (milliers de \$)},tick label style={font=\footnotesize},
  grid=major,grid style={gray!20},table/col sep=comma]
\addplot[only marks,mark=-,omDef,thick,error bars/.cd,y dir=both,y explicit,error bar style={line width=5pt,omDef!45}]
  table[x=grade,y=mid,y error=half]{figdata/industry/09-exchanges-brokers-vendors-and-regulators-as-employers/skbands.csv};
\draw[omThm,dashed] (axis cs:8.4,292.3) -- (axis cs:17.6,292.3);
\node[font=\scriptsize,omThm,anchor=south west] at (axis cs:8.5,292.3) {plafond 292\,300~\$};
\end{axis}
\end{tikzpicture}

\omcaption{\omterm{def:in:exchanges-brokers-vendors-and-regulators-as-employers:band}{Bandes de rémunération} 2026 de la SEC pour les échelons SK-9 à SK-17 à New York~: salaire de base multiplié par un plus le taux de majoration locale de 37.95~\%,
plafonné à 292\,300~\$ (tirets). Chaque barre va du minimum de l'échelon à son maximum. Données~: SEC, Compensation overview,
par l'intermédiaire de \texttt{firm.payband.at}.}\label{fig:in:exchanges-brokers-vendors-and-regulators-as-employers:bands}
\end{omfigure}

\section{Ce que l'infrastructure propose dans les dépôts}

Les quatre groupes de bourses déposent aussi des dossiers de conditions de travail (chapitre~14), et leurs offres peuvent être
rapprochées de celles des banques et des teneurs de marché, ainsi que de la bande du régulateur.

\begin{dated}{2025-09}{Groupes de bourses dans les dépôts}\label{dat:in:exchanges-brokers-vendors-and-regulators-as-employers:lca}
Exercice fiscal 2025, dossiers de conditions de travail, salaire de base médian proposé~: groupes de bourses, tous intitulés 130\,800~\$ (462 dossiers,
4 employeurs), logiciel 140\,400~\$ (329), recherche quantitative 107\,100~\$ (35), apprentissage automatique et données 108\,500~\$ (60)~;
banques 154\,260~\$, 155\,688~\$, 158\,100~\$ et 140\,714~\$~; teneurs de marché 175\,000~\$ dans chacun des quatre cas. Salaire de base uniquement.
\end{dated}

\begin{omfigure}
\begin{tikzpicture}
\begin{axis}[ybar,area legend,width=10.5cm,height=5.4cm,xmin=0.5,xmax=4.5,ymin=0,ymax=200,bar width=10pt,xtick={1,2,3,4},
  xticklabels={{tous postes},{logiciel},{recherche},{ML et données}},
  ylabel={\small salaire de base médian proposé, milliers de \$},tick label style={font=\footnotesize},grid=major,grid style={gray!20},
  table/col sep=comma,legend style={legend cell align=left,font=\scriptsize,at={(0.5,-0.25)},anchor=north,legend columns=3,
  /tikz/every even column/.append style={column sep=8pt}}]
\addplot[fill=black!30,draw=none] table[x=pos,y=exchange]{figdata/industry/09-exchanges-brokers-vendors-and-regulators-as-employers/filings.csv};
\addplot[fill=omDef,draw=none] table[x=pos,y=bank]{figdata/industry/09-exchanges-brokers-vendors-and-regulators-as-employers/filings.csv};
\addplot[fill=omThm!70,draw=none] table[x=pos,y=market_maker]{figdata/industry/09-exchanges-brokers-vendors-and-regulators-as-employers/filings.csv};
\draw[black!60,dashed] (axis cs:0.5,157.698) -- (axis cs:4.5,157.698);
\legend{groupes de bourses,banques,teneurs de marché}
\end{axis}
\end{tikzpicture}

\omcaption{Salaire de base médian proposé dans les dossiers de conditions de travail de l'exercice fiscal 2025, par famille de métiers et type d'employeur, avec
le bas de la bande SK-14 de la SEC à New York (tirets). Données~: \texttt{data/industry/lca\_ranges.csv}, par
\texttt{in\_infra.filings}.}\label{fig:in:exchanges-brokers-vendors-and-regulators-as-employers:filings}
\end{omfigure}

Les offres des bourses se situent sous celles des banques dans chaque famille, et nettement sous celles des teneurs de marché
(\cref{fig:in:exchanges-brokers-vendors-and-regulators-as-employers:filings})~: 130\,800~\$ contre 154\,260~\$ et
175\,000~\$ pour l'ensemble des intitulés~; pour la recherche, 107\,100~\$ contre 158\,100~\$ et 175\,000~\$. La bande du régulateur se situe entre les deux~:
un poste SK-14 à New York commence à 157\,698~\$ et va jusqu'à 267\,168~\$, au-dessus des offres médianes des bourses et dans
la plage des banques. Aucun de ces chiffres n'inclut la prime, plus élevée dans les sociétés de trading (chapitre~14)~; la bande publique
constitue toute la rémunération du régulateur.

\section{Passer d'un côté à l'autre}

Les professionnels circulent entre l'industrie et son infrastructure dans les deux sens~: des ingénieurs passent des bourses aux sociétés de trading
et inversement, des quants passent des régulateurs aux banques et des banques aux régulateurs, des économistes de la structure de marché des plateformes
de négociation vers les régulateurs. Les départs d'un régulateur sont soumis aux règles post-emploi de la fonction publique, et les arrivées dans un régulateur relèvent des règles sur les
conflits d'intérêts~; les unes comme les autres sont les règles de l'employeur, fixées dans ses propres codes. Pour une carrière, l'infrastructure est
l'endroit où l'on apprend comment fonctionne le marché de l'intérieur --- sa plomberie, ses règles et ses données --- avec une rémunération
publiée, plus stable, et en général plus basse au sommet que celle des sociétés de trading (chapitre~14).

\section{Tutoriel~: revenu par tête dans l'infrastructure}\label{tut:in:exchanges-brokers-vendors-and-regulators-as-employers:tut}

\begin{tutorial}
\textbf{Objectif.} Calculer, à partir de leurs dépôts, le revenu et le résultat opérationnel par tête des groupes de bourses, et
convertir la grille publiée par un régulateur en plages annuelles. \textbf{État final~:}
les \cref{fig:in:exchanges-brokers-vendors-and-regulators-as-employers:perhead,fig:in:exchanges-brokers-vendors-and-regulators-as-employers:bands}.
\begin{tutsteps}
\item \textbf{Les dépôts.} \texttt{data/industry/infrastructure\_employers.csv} contient pour chaque groupe le revenu total et
le résultat opérationnel (tiré des données SEC XBRL company facts lorsqu'elles sont balisées, sinon du compte de résultat) et son effectif de fin d'exercice (tiré
du texte du 10-K)~; \texttt{in\_infra.exchanges()} calcule les mesures par tête et la marge opérationnelle.
\item \textbf{La grille.} \texttt{firm.payband.load\_scale} lit les échelons SK, les taux de majoration locale et le plafond~;
\texttt{firm.payband.at(scale, grade, station)} renvoie la plage annuelle
(le \cref{lst:in:exchanges-brokers-vendors-and-regulators-as-employers:at}).
\omcode{code/firm/payband/firm_payband.py}{41}{52}{Plage annuelle d'un échelon dans un lieu d'affectation, et recouvrement de deux
plages.}\label{lst:in:exchanges-brokers-vendors-and-regulators-as-employers:at}
\item \textbf{Comparer.} \texttt{firm.payband.overlap(a, b)} donne la part d'une plage $b$ couverte par une bande $a$, pour
la comparaison du chapitre~14 avec les dossiers de conditions de travail.
\end{tutsteps}
\end{tutorial}

\textbf{Pour aller plus loin.} Calculer le revenu par tête de chaque groupe, net des frais liés aux transactions lorsque le 10-K
les indique (exercice~7)~; dresser le tableau des bandes SK pour Chicago et Washington, et les comparer.

\section{Construction~: grilles de rémunération publiques}\label{bld:in:exchanges-brokers-vendors-and-regulators-as-employers:payband}

\begin{build}
\bfield{Rôle} Lire les grilles de rémunération publiées du secteur public comme des plages annuelles, afin que le chapitre~14 puisse les placer à côté des
plages du secteur privé issues des dossiers de conditions de travail, et le chapitre~24 à côté de celle des fonctions de contrôle.
\bfield{Interface} \texttt{firm.payband}~: \texttt{Band(grade, lo, hi)}, \texttt{Scale(name, bands, locality, cap)}~;
\texttt{load\_scale(name, bands\_csv, locality\_csv, cap)}~; \texttt{at(scale, grade, station)}~; \texttt{overlap(a, b)}.
\bfield{Règles} La majoration locale s'applique au salaire de base~; le plafond s'applique au total, à chaque extrémité de la plage~; un lieu
d'affectation ou un échelon inconnu est une erreur.
\bfield{Tests d'acceptation} \texttt{code/firm/payband/tests/}~: majoration locale et plafond sur une grille construite~; recouvrement de
plages disjointes et qui se chevauchent~; chargement depuis un CSV.
\bfield{Pour aller plus loin} Tables par paliers au sein d'un échelon (rémunération selon l'ancienneté)~; plusieurs grilles dans une même monnaie pour une comparaison
internationale (chapitre~27).
\end{build}

\begin{omsources}
\item Formulaires 10-K pour 2025 de CME Group, Intercontinental Exchange, Nasdaq et Cboe Global Markets~; SEC XBRL company facts.
\item SEC, Compensation overview (grille SK 2026 et taux de majoration locale).
\item CPMI-IOSCO, Principles for financial market infrastructures (2012), page de publication de la BIS.
\end{omsources}

\section{Exercices}

\begin{exercise}[$\star$]\label{exo:in:exchanges-brokers-vendors-and-regulators-as-employers:1}
Calculer la marge opérationnelle de chaque groupe boursier en 2025.
\end{exercise}

\begin{exercise}[$\star$]\label{exo:in:exchanges-brokers-vendors-and-regulators-as-employers:2}
Calculer la fourchette de New York d'une position SK-13, ainsi que le rapport entre son haut et son bas.
\end{exercise}

\begin{exercise}[$\star$]\label{exo:in:exchanges-brokers-vendors-and-regulators-as-employers:3}
Pour quels grades la borne supérieure de New York est-elle fixée par le plafond plutôt que par le grade~?
\end{exercise}

\begin{exercise}[$\star\star$]\label{exo:in:exchanges-brokers-vendors-and-regulators-as-employers:4}
Pourquoi le revenu par tête est-il une moins bonne base de comparaison entre les quatre groupes boursiers que le résultat opérationnel par tête~?
\end{exercise}

\begin{exercise}[$\star\star$]\label{exo:in:exchanges-brokers-vendors-and-regulators-as-employers:5}
Une position SK-14 à Chicago et une position à New York~: calculer les deux fourchettes ainsi que l'écart entre leurs milieux.
\end{exercise}

\begin{exercise}[$\star\star$]\label{exo:in:exchanges-brokers-vendors-and-regulators-as-employers:6}
À partir de la \cref{fig:in:exchanges-brokers-vendors-and-regulators-as-employers:perhead}, quel groupe présente l'écart le plus large
entre revenu et résultat opérationnel par tête, et qu'est-ce qui pourrait l'expliquer~?
\end{exercise}

\begin{exercise}[$\star\star\star$]\label{exo:in:exchanges-brokers-vendors-and-regulators-as-employers:7}
\emph{Programmation.} Nasdaq annonce des revenus totaux diminués des frais liés aux transactions de 5\,249~millions de dollars. Recalculer son revenu par
tête sur cette base et sa marge opérationnelle sur cette même base. Quelle part de l'écart de marge opérationnelle avec la bourse de produits dérivés se résorbe~?
\end{exercise}

\begin{exercise}[$\star\star\star$]\label{exo:in:exchanges-brokers-vendors-and-regulators-as-employers:8}
\emph{Trouver la faille.} «~Un examinateur SK-16 à New York gagne 292\,300~\$, et un quant dans une société de trading gagne moins que cela
en salaire de base, donc le régulateur paie davantage.~»
\end{exercise}

\section{Problème~: revenu par tête dans l'infrastructure}

\begin{problem}[{Problème du week-end --- revenu par tête dans l'infrastructure}]\label{pb:in:exchanges-brokers-vendors-and-regulators-as-employers:1}
Un ingénieur ayant reçu des offres d'un groupe boursier et d'un régulateur veut comprendre les deux employeurs à partir de leurs
chiffres publics.

\medskip\noindent\textbf{Partie I --- L'infrastructure.}
\begin{enumerate}
\item Définir une \omterm{def:in:exchanges-brokers-vendors-and-regulators-as-employers:fmi}{infrastructure de marché financier} et en lister les cinq types.
\item Que vend un groupe boursier, et quels emplois de quant et d'ingénierie offre-t-il~?
\item Donner le nombre de salariés des quatre groupes en 2025.
\item Donner leur revenu et leur résultat opérationnel par tête.
\item Qu'est-ce qui gonfle le revenu de trois d'entre eux, et de quelle manière~?
\end{enumerate}

\medskip\noindent\textbf{Partie II --- La comparaison.}
\begin{enumerate}[resume]
\item Donner les quatre marges opérationnelles.
\item Combien de fois le résultat opérationnel par tête le plus élevé est-il supérieur au plus faible~?
\item Quels segments hors négociation deux des groupes publient-ils~?
\item Que vend un courtier, et que vend un fournisseur de données~?
\item Qu'est-ce qui distingue la rémunération d'un fournisseur de technologies de trading de celle d'une société de trading~?
\end{enumerate}

\medskip\noindent\textbf{Partie III --- Le régulateur.}
\begin{enumerate}[resume]
\item Définir une \omterm{def:in:exchanges-brokers-vendors-and-regulators-as-employers:band}{bande de rémunération}.
\item Énoncer la méthode qui transforme une bande en fourchette annuelle.
\item Donner les fourchettes de New York des grades SK-12, SK-14 et SK-16.
\item Quels grades sont plafonnés à New York, et à quel niveau~?
\item Pourquoi les bandes sont-elles si larges, et où un nouveau recruté est-il placé~?
\end{enumerate}

\medskip\noindent\textbf{Partie IV --- Le verdict.}
\begin{enumerate}[resume]
\item Énoncer le \emph{résultat nommé}~: la fourchette du résultat opérationnel par tête des quatre groupes, et
la fourchette de New York d'une position SK-14.
\item Quelles règles régissent les départs vers un régulateur et les arrivées depuis un régulateur~?
\item Qu'est-ce que l'infrastructure enseigne qu'une société de trading n'enseigne pas~?
\item Où le chapitre~14 compare-t-il ces bandes à la rémunération privée~?
\item En deux phrases, comment l'ingénieur doit-il comparer les deux offres~?
\end{enumerate}
\end{problem}

\section{Questions d'entretien}

\begin{interviewq}[$\star$ \iqroles{développeur}]\label{iq:in:exchanges-brokers-vendors-and-regulators-as-employers:1}
Qu'est-ce qui préoccupe l'équipe du moteur d'appariement (matching engine) d'une bourse, et qui ne préoccupe pas les ingénieurs de stratégie d'une société de trading~?
\end{interviewq}

\begin{interviewq}[$\star$ \iqroles{risque}]\label{iq:in:exchanges-brokers-vendors-and-regulators-as-employers:2}
Que doit faire le modèle de marge d'une chambre de compensation, et qui supporte le risque si celui-ci est erroné~?
\end{interviewq}

\begin{interviewq}[$\star\star$ \iqroles{développeur}]\label{iq:in:exchanges-brokers-vendors-and-regulators-as-employers:3}
Une bourse veut que chaque participant d'une salle de colocation reçoive les données de marché au même instant. Comment concevriez-vous
et testeriez-vous cela~?
\end{interviewq}

\begin{interviewq}[$\star\star$ \iqroles{chercheur}]\label{iq:in:exchanges-brokers-vendors-and-regulators-as-employers:4}
Un régulateur vous demande d'estimer la part du volume d'un marché qui relève des courses à la latence. Quelles données vous faut-il, et quelle est
la première grandeur à mesurer~?
\end{interviewq}

\begin{interviewq}[$\star\star$ \iqroles{chercheur}]\label{iq:in:exchanges-brokers-vendors-and-regulators-as-employers:5}
Un fournisseur de données évalue une obligation d'entreprise qui n'a pas été échangée depuis trois semaines. Comment s'y prend-il, et comment savez-vous que le prix est juste~?
\end{interviewq}

\begin{interviewq}[$\star\star\star$ \iqroles{risque, développeur}]\label{iq:in:exchanges-brokers-vendors-and-regulators-as-employers:6}
Une bourse verse des rabais aux fournisseurs de liquidité et les comptabilise en frais liés aux transactions. Comment ces rabais modifient-ils la
signification de son chiffre d'affaires, et comment compareriez-vous équitablement deux bourses~?
\end{interviewq}
